%% maths-cercle-quadrants — les quatre quadrants du cercle trigonométrique.
%% Chaque quart de disque est légèrement teinté et numéroté I, II, III, IV dans le sens
%% direct (antihoraire), rappelé par la flèche courbe « + » autour de O ; l'abscisse d'un point
%% du cercle est son cosinus, son ordonnée son sinus (axes annotés).
%% Les quatre étiquettes de signes (cos > 0 / sin < 0 …) sont enveloppées de \rappel :
%% la variante « maths-cercle-quadrants-muet » ne montre que les quadrants nus, pour
%% demander à l'élève dans quel quadrant tombe un angle donné, et quels signes en découlent.
\documentclass[tikz,border=4pt]{standalone}
\usepackage{liaisons}
\begin{document}
\begin{tikzpicture}[x=2cm,y=2cm,
  axe/.style={-{Stealth[length=5pt]}, line width=0.6pt},
  sens/.style={-{Stealth[length=6pt]}, line width=0.8pt},
  num/.style={font=\large, text=black!55},
  signe/.style={font=\scriptsize, align=center, inner sep=3pt, line width=0.5pt,
                fill=white, rounded corners=2pt}]

%% ---------------- les quatre quarts de disque teintés ------------------
%% \angle de départ / couleur / numéro
\foreach \dep/\coul/\num in {0/solideE/I, 90/solideB/II, 180/solideF/III, 270/solideD/IV} {
  \fill[\coul!12] (0,0) -- (\dep:1) arc (\dep:\dep+90:1) -- cycle;
  \node[num] at (\dep+45:0.72) {\num};}

%% ---------------- repère et cercle -------------------------------------
\draw[axe] (-1.18,0) -- (1.22,0)
     node[anchor=south east, inner sep=3pt, font=\small, text=black!60] {$\cos$};
\draw[axe] (0,-1.18) -- (0,1.22)
     node[anchor=north west, inner sep=3pt, font=\small, text=black!60] {$\sin$};
\draw[line width=1.1pt, black!75] (0,0) circle (1);
\node[font=\small, anchor=north east, inner sep=1pt, fill=white] at (-0.02,-0.02) {$O$};
\fill (1,0) circle (1.7pt);
\node[font=\small, anchor=north east, inner sep=3pt] at (1,0) {$I$};
\fill (0,1) circle (1.7pt);
\node[font=\small, anchor=north east, inner sep=3pt] at (0,1) {$J$};

%% ---------------- sens direct (flèche courbe autour de l'origine) ------
\draw[sens] (8:0.40) arc (8:102:0.40);
\node[font=\small, inner sep=1pt] at (-16:0.44) {$+$};

%% ---------------- signes de cos et sin, quadrant par quadrant ----------
%% C'est la réponse : effacé dans la variante muette.
\rappel{
\node[signe, draw=solideE, anchor=south west] at (0.86,0.86)
     {$\cos>0$\\ $\sin>0$};
\node[signe, draw=solideB, anchor=south east] at (-0.86,0.86)
     {$\cos<0$\\ $\sin>0$};
\node[signe, draw=solideF, anchor=north east] at (-0.86,-0.86)
     {$\cos<0$\\ $\sin<0$};
\node[signe, draw=solideD, anchor=north west] at (0.86,-0.86)
     {$\cos>0$\\ $\sin<0$};
}
\end{tikzpicture}
\end{document}
